\begin{frame}{대칭을 깨뜨리면 $b$만 변한다 (병진 예제)}
  실제 데이터는 원점에 대칭인 경우가 거의 없다.
  \begin{itemize}
    \item 같은 6개 데이터를 전부 $+2$ 오른쪽으로 병진하면
      양성은 $(5,3),(6,3),(5,4)$, 음성은 $(-1,-3),(-2,-3),(-1,-4)$.
    \item 최적해: $w=(1/6,1/6)$ — \textbf{변함없음},
      그런데 $b = -1/3$.
    \item 왜 $-1/3$인가: 병진 후 원점 방향으로 가장 가까운 양의 점
      $(5,3)$을 마진 경계에 놓으면:
      \[
      1\times\Bigl(\tfrac16\cdot5+\tfrac16\cdot3+b\Bigr)
      = \tfrac86 + b = 1
      \;\Longrightarrow\; b = -\tfrac13
      \]
    \item 음성 쪽 가장 가까운 점 $(-1,-3)$도 같은 $b$로 등호 확인 —
      새로운 두 서포트 벡터: $(5,3)$과 $(-1,-3)$.
  \end{itemize}
  \vskip0.4em
  \begin{block}{본질}
    \textbf{데이터 전체를 이동시켜도 $w$는 그대로, $b$만 변한다} —
    ``마진 최대화''의 본질은 두 클래스 간 \textbf{상대적 방향}($w$)을
    정하는 것이며, \textbf{어디에 놓이느냐}($b$)는 데이터의 위치에
    따라 달라진다.
  \end{block}
\end{frame}
